\section{Análisis de rendimiento}

Para caracterizar el rendimiento del sistema se han medido tres métricas principales:
uso de CPU del proceso voctocore, consumo de memoria RAM y latencia extremo a extremo
(desde la captura de un fotograma en la fuente hasta su visualización en el monitor
principal de voctogui).

Las mediciones se realizaron en el Escenario 1 (monopuesto) con el sistema operando
con cuatro fuentes FFmpeg activas simultáneamente a 1920×1080 @25\,fps, durante
sesiones de prueba de 30 minutos.

\subsection{Uso de CPU}

El proceso voctocore consume la mayor parte de los recursos de CPU, ya que concentra
toda la composición de vídeo mediante el pipeline GStreamer. El uso observado varía
en función del número de fuentes activas y del modo de composite empleado.

\subsection{Latencia}

La latencia extremo a extremo en el Escenario 1 (loopback TCP) es inferior a la de
un fotograma a 25\,fps ($<$40\,ms), lo que garantiza una experiencia de realización
fluida sin retardos perceptibles entre la acción del operador y el efecto en el
monitor de programa.

En el Escenario 2 (red local 1\,GbE), la latencia aumenta en función de la
configuración de jitter buffer de GStreamer, manteniéndose en valores adecuados para
producción en directo.

% --- Tabla sugerida ---
% \begin{table}[H]
%   \centering
%   \caption{Métricas de rendimiento por escenario (4 fuentes 1080p @25fps).}
%   \label{tab:rendimiento}
%   \begin{tabular}{|l|c|c|c|}
%     \hline
%     \textbf{Escenario} & \textbf{CPU voctocore (\%)} & \textbf{RAM (MB)} & \textbf{Latencia (ms)} \\
%     \hline
%     1 --- Monopuesto    & xx\%  & xxx & $<$40 \\
%     2 --- Red local     & xx\%  & xxx & xx--xx \\
%     3 --- Docker        & xx\%  & xxx & xx--xx \\
%     \hline
%   \end{tabular}
% \end{table}

% --- Figura sugerida ---
% \begin{figure}[H]
%   \centering
%   \includegraphics[width=0.7\textwidth]{figuras/cpu_vs_fuentes.png}
%   \caption{Uso de CPU de voctocore en función del número de fuentes activas.}
%   \label{fig:cpu_vs_fuentes}
% \end{figure}
